\section{Inference Optimizations}

\subsection{Model Acceleration}
\textbf{DiT Optimizations.} To accelerate DiT inference, we adopt diffusion distillation techniques to reduce the number of function evaluations (NFE) required for generation. We incorporate the Trajectory Segmented Consistency Distillation (TSCD) technique, originally introduced in HyperSD\cite{ren2025hyper}, which partitions the denoising trajectory into multiple segments and enforces consistency between predicted and target states across these segments. This allows the student model to learn an accurate approximation of the diffusion process with fewer steps. Using TSCD, our DiT model performs competitively with 4x acceleration, offering a strong balance between speed and fidelity. To push acceleration further, we incorporate Score Distillation from  RayFlow\cite{shao2025rayflow}, which aligns the student model’s predicted noise (i.e., score function) with that of the teacher using expected noise consistency. This approach supports trajectory-level optimization for each sample, enabling more stable and adaptive sampling even at low NFEs. It effectively improves generalization and reduces artifacts during fast generation. To improve visual quality, we extend the adversarial training strategy from APT\cite{lin2025diffusion} to a multi-step distillation setting, incorporating human preference data for supervision. A learned discriminator guides the student model toward outputs favored by human judgments, effectively mitigating artifacts from aggressive acceleration and enhancing perceptual realism.

Through the proposed distillation pipeline, our final distilled model achieves comparable results to the original model across four expert-evaluated dimensions: prompt alignment, motion quality, visual fidelity, and consistency with the source image.

\textbf{VAE Optimizations.} In video generation tasks, the decoding process from latent space to pixel space incurs significant computational cost. We profiled the VAE decoder and found that stages closer to the pixel space dominate the latency. By narrowing the channel widths in these stages, we design a thin VAE decoder. Retraining it with a fixed pre-trained encoder, we achieve a 2× speedup with no loss in visual quality of the end-to-end video generation.

\subsection{Inference Infrastructure}

\textbf{High-Performance Kernel.} Extensive kernel fusion efforts have been conducted on the model's core modules, resulting in a cumulative 15\% improvement in the model's inference throughput.

\textbf{Quantization and Sparse.} Building on the Seedream \cite{gong2025seedream} technical solution, we have implemented fine-grained mixed-precision quantization tailored for Attention and Gemm operations. Moreover, our exploration revealed that the sparse attributes of DiTs exhibit hierarchical and blockified structures across and within various modalities. Expanding on the methodology established by AdaSpa \cite{adaspa}, we have introduced a streamlined tuning solution focused on minimizing search stage overhead. Additionally, we have successfully integrated our optimized fine-grained Attention Quantization approach into this scheme. Numerous efforts have been dedicated to mitigating the impact of full quantization and sparsity on the quality of pixel-level generation. We have achieved an optimal balance between performance and efficiency.

\textbf{Parallelism Strategy.} In order to decrease the allocated massive memory due to the long sequence in video generation schema. A customized adaptive hybrid parallel strategy has been proposed to effectively split the sequences. This approach integrates the concept of context parallelism to optimize communication processes, resulting in a reduction of communication overhead to a quarter of the level observed in Ulysses \cite{jacobs2023deepspeed}. Simultaneously, we have further reduced end-to-end communication overhead by introducing FP8 communication. 

\textbf{Async Offloading Strategy.} Due to the extensive computational demands of attention coupled with the large model size. We developed an automated and adaptive AsyncOffloading strategy. We successfully solved the large model deployment problem on memory-limited devices with a performance drop of less than 2\%.

\textbf{Hybrid Parallelism for Distributed VAE.} Moreover, to address the issue of high GPU memory consumption due to the VAE-Decoder, we implemented an adaptive hybrid parallel strategy. This method partitions the input data along the spatial and temporal dimensions simultaneously and employs efficient collective communication for Conv3D computation. Thus, we further improved parallel scaling performance.

\textbf{Pipeline Optimizations.} We adopted kernel fusion, quantization, parallelization, continuous batching, prefix caching, and other common techniques to improve the overall throughput of the prompt engineering effectively. Furthermore, to tackle the issue of low encoding efficiency in long videos, we have implemented video encoding acceleration solutions. 

These innovations have effectively boosted the E2E efficiency of the whole inference pipeline.